\documentclass[11pt]{article}
\usepackage[margin=1in]{geometry}
\usepackage{amsmath,amssymb}
\usepackage{booktabs}
\usepackage{graphicx}
\usepackage{hyperref}
\usepackage{natbib}

\title{Policy-Guided Information-Set Search in Four-Player Imperfect
Information Card Games: An Empirical Study in Big~2}

\author{%
  Brandon Luo\\
  \texttt{bluo42/big2-ai}
  \and
  Claude (Anthropic)\\
  \texttt{engineering and analysis}
}

\date{August 2026}

\begin{document}
\maketitle

\begin{abstract}
We report an empirical study of combining proximal policy optimization
(PPO) with information-set Monte Carlo tree search (IS-MCTS) in Big~2,
a four-player imperfect-information climbing card game with a
$\sim 10^{18}$ information-set space and no partnership structure. Over
a sequence of controlled experiments we ablate nine engineered feature
blocks against a paired evaluation protocol, measure where search
actually changes decisions, and evaluate the resulting agents against
each other and against 570 recorded games from human players. Four
findings stand out. First, \emph{search-assisted rollouts during
training did not pay}: in six paired comparisons across two independent
lineages, training with the tree in the loop cost $0.03$--$0.40$
points/game relative to policy-only training at $8$--$12\times$ the
wall-clock. Second, once a policy is trained on a sufficiently diverse
opponent diet, \emph{IS-MCTS and the policy converge}: the search
overrides the policy on only $0.4\%$ of in-distribution decisions, a
figure invariant to a $32\times$ increase in simulations, a $3.5\times$
reduction in override margin, deeper playouts, and horizon changes ---
but it doubles against out-of-distribution opponents. Third,
\emph{evaluation noise in four-player card games is severe enough to
invert conclusions}: single-seed 600-game evaluations of the same fixed
checkpoint spanned $-0.03$ to $+2.09$ points/game, and one feature block
was rejected and then reinstated after re-adjudication under a paired
three-seed protocol. Fourth, the trained agents are \emph{strongly
non-transitive}: a 35-composition round-robin over $35{,}000$ games
exhibits a rock-paper-scissors cycle among three agents whose pairwise
margins reverse depending on table composition. We close with the
implications for bridge and mahjong.
\end{abstract}

\section{Introduction}

Perfect-information game search has a mature theory; imperfect
information does not, and the four-player non-partnership setting is
particularly awkward. Determinized search (PIMC) is known to suffer
strategy fusion and non-locality \citep{frank1998,long2010}, yet it
remains competitive in practice. Deep RL systems for card games ---
DouZero \citep{zha2021}, PerfectDou \citep{yang2022}, Suphx
\citep{li2020} --- typically train a policy or value network with
self-play and either omit search at play time or add it late. Which
component actually carries the strength, and where, is rarely measured
directly.

This paper is an engineering-grade empirical answer to that question in
one domain. Big~2 (\emph{Deuces}, \begin{CJK}{UTF8}{gbsn}锄大地\end{CJK})
is a four-player climbing game: 52 cards dealt evenly, players shed
singles, pairs, and five-card poker hands in ascending order, and the
first player out collects the others' remaining cards under a tiered
multiplier. It is imperfect-information, non-partnership, has a
combinatorially awkward action space ($\sim 10^2$ legal moves in the
opening), and --- crucially for this study --- collapses into an
exactly solvable perfect-information-like endgame once few enough cards
remain.

Our contributions are empirical rather than algorithmic:

\begin{enumerate}
\item A \textbf{feature-ablation ladder} that adds engineered blocks one
  at a time under a paired gate, showing that endgame-danger features
  pay on every lineage tested, planning/dominance features pay on one
  lineage and hurt on another, and cross-game opponent inference pays
  nowhere.
\item A \textbf{negative result on search-in-training}, consistent in
  sign across six paired comparisons.
\item A \textbf{decision census} localizing where search and exact
  solving actually change play.
\item A \textbf{non-transitivity measurement} over all 35 four-seat
  compositions of four trained agents.
\item A \textbf{human-deviation study} over 570 recorded games
  identifying the residual leak as early-hand patience, not endgame
  tactics.
\end{enumerate}

\section{Testbed and evaluation}

\paragraph{Rules.} Ranks $3 < 4 < \cdots < K < A < 2$; suits
$\diamondsuit < \clubsuit < \heartsuit < \spadesuit$. The holder of
$3\diamondsuit$ opens and must include it. Legal plays are singles,
pairs, and five-card hands (straight $<$ flush $<$ full house $<$ quads
$<$ straight flush); lone triples are disabled. Passing is soft: a pass
skips the turn but the player may re-enter when the trick returns.
Payout: losers pay the winner their remaining card count, doubled at
$\geq 10$ cards and tripled at 13.

\paragraph{Scoring metric.} All results are mean points/game from a
rotated seat. Because Big~2 is zero-sum across four seats, a score of
$0$ means parity with the field, and $\pm 1$ point/game is a large
margin (a typical decisive game swings $5$--$15$ points).

\paragraph{The evaluation-variance problem.} Our initial gate used one
seed and 600 games. A variance audit of two \emph{fixed} checkpoints
across three seeds produced the scores in Table~\ref{tab:variance}.

\begin{table}[h]
\centering
\begin{tabular}{lccc}
\toprule
Checkpoint & seed 11 & seed 22 & seed 33 \\
\midrule
A/plan (600 games) & $-0.033$ & $+1.777$ & $+2.087$ \\
A/beatprob (600 games) & $+0.083$ & $+1.748$ & $+2.277$ \\
A/plan (1500 games) & $+0.340$ & $+1.481$ & $+1.241$ \\
A/beatprob (1500 games) & $+0.434$ & $+1.423$ & $+1.610$ \\
\bottomrule
\end{tabular}
\caption{The same fixed checkpoints re-evaluated across seeds. Absolute
scores swing by more than 2 points/game; the seed drives both deal
sequence and opponent draw, so it is a table-composition effect, not
merely deal luck.}
\label{tab:variance}
\end{table}

The correct protocol is therefore \emph{paired}: identical seeds for
both checkpoints, differences averaged across seeds. Under
$3 \times 1500$ paired games the difference standard deviation falls to
$\approx 0.2$ points/game. One feature block (\S\ref{sec:ladder}) was
rejected under the single-seed gate and reinstated under the paired
one; we regard single-seed evaluation in this domain as unusable.

\section{Agent architecture}

The policy is a set-attention pointer network: the state and each legal
move are embedded by two-layer MLPs to width $d$, summed, passed
through masked multi-head self-attention over the candidate-move set,
and scored by a linear head. Three heads share the trunk: a policy
pointer over moves, a value head, and a 156-dimensional belief head
(per-opponent card posterior) trained as an auxiliary loss.

\paragraph{Attention depth.} All shipped models use two attention
blocks. The second block is added to an existing checkpoint with a
zero-initialized output projection under a pre-norm residual, making
the widened network \emph{bit-identical} to its predecessor at
initialization (verified: max logit difference $0.00$). After
$\sim 160$k games of training the second block's output projection had
grown to weight norm $53.5$ versus $11.9$ for the original block,
indicating that gradient flow moved the bulk of move-interaction
processing into the new block.

\paragraph{Feature blocks.} Move-side features, in the order they were
introduced historically:
\begin{itemize}
\item \textbf{base} (296 dims): DouZero-style action encoding, analytic
  beliefs from deterministic card exclusion, hand-structure and
  decomposition proxies, payment-tier risk, plus a linear champion's
  feature vector and scalar advice.
\item \textbf{danger} ($+14$): thresholded opponent card counts
  (exactly 1, 2, $\leq 3$), which seat is nearly out, cheap-single-gift
  detection, can-anyone-go-out-on-this-combo-size.
\item \textbf{plan} ($+16$ move, $+24$ state): exact boss detection (is
  this move unbeatable given the unseen set --- the ``free turn''),
  guaranteed-win, keeps-the-lead, wastes-control, race margin, denial
  vs.\ gift, share of unseen mass above the move; plus run-out summary
  and opponent-read state features.
\item \textbf{profiles} ($+24$ state): cross-game opponent style
  estimates on a 500-game EMA.
\item \textbf{beat-risk} ($+4$): per candidate move, $P(\text{next
  opponent can answer})$, $P(\text{any opponent can answer})$,
  $E[\#\text{answerers}]$, and $P(\text{any opponent holds the
  class})$, computed by uniform Monte Carlo over consistent deals of
  the unseen pool. This is deliberately \emph{low-bias}: deterministic
  exclusion of played cards plus a uniform prior, no learned or
  play-evidence inference. It makes the flush-versus-straight decision
  directly visible to the network.
\end{itemize}

\section{Play-time search}

The deployed agent is a single decision procedure with no stage switch:

\begin{enumerate}
\item \textbf{Forced moves} are played immediately.
\item \textbf{Exact solving.} If a move empties the hand and is
  unbeatable given the unseen set, it is played on the spot. Below 14
  total cards, every candidate is evaluated by max$^n$ search to
  terminal states over belief-sampled worlds (weighted PIMC), with a
  node budget, a wall-clock deadline, and an agreement threshold; a
  process-level transposition cache retains only \emph{exactly} solved
  positions, so repeated endings are answered from the table.
\item \textbf{IS-MCTS} elsewhere: PUCT selection with the policy as
  prior, candidate set pruned to the moves covering $90\%$ of prior
  mass (at most 6), determinized worlds drawn per simulation,
  opponents' in-tree replies sampled from the agent's own policy,
  playouts truncated at the player's next decision --- or earlier, the
  moment a trick resolves with the player holding the fresh lead, since
  that is the highest-leverage state to price --- and leaves evaluated
  by the value head or solved exactly when in range.
\item \textbf{Margin.} The search's move is played only if its measured
  advantage over the policy's move exceeds a threshold.
\end{enumerate}

\section{Experiments}

\subsection{The feature ladder}
\label{sec:ladder}

Two lineages were rebuilt from scratch at $d=192$ with two attention
blocks: lineage~A warm-started from the original PPO champion,
lineage~B from a trained exploiter (an agent trained as a best response
to the champion). Each stage trains $120$k games with one additional
feature block, zero-initialized at the input layer, then gates on the
paired champions-room score. Table~\ref{tab:ladder} reports the
outcome.

\begin{table}[h]
\centering
\begin{tabular}{lcc}
\toprule
Block & Lineage A & Lineage B \\
\midrule
base (diet only) & $+0.985$ abs. & $+1.217$ abs. \\
danger & \textbf{kept} $+0.587$ & \textbf{kept} $+0.148$ \\
plan & \textbf{kept} $+0.158$ & \textbf{dropped} $-0.193$ \\
profiles & removed by design & removed by design \\
beat-risk & \textbf{kept} $+0.135^{\dagger}$ & not applicable \\
search-in-training & \textbf{dropped} $-0.288^{\dagger}$ & \textbf{dropped} $-0.154^{\dagger}$ \\
\bottomrule
\end{tabular}
\caption{Feature-ladder verdicts, gains in points/game.
$^{\dagger}$paired $3\times1500$ protocol; earlier rows used the
single-seed 600-game gate and should be read as indicative only. The
beat-risk block was rejected by the single-seed gate ($-0.160$) and
reinstated by the paired one ($+0.135$).}
\label{tab:ladder}
\end{table}

Three observations. (i) The \emph{diet} mattered more than any feature:
simply training against a weighted draw from the full model collection
(all scripted baselines, linear and evolutionary champions, both
shipped networks, a humanlike clone, and a dedicated exploiter, with
the models-to-beat at double weight) moved lineage~A from $+0.26$ to
$+0.99$ before any new feature was added. (ii) Endgame-danger features
pay on both lineages. (iii) The planning block's value is
\emph{lineage-dependent}: kept on A, harmful on B, which we attribute
to the exploiter lineage having already internalized dominance
reasoning in its weights.

\paragraph{Opponent profiles.} We removed the cross-game opponent-style
block on principle (it is inference over opponent identity, not game
state) and then discovered it had been \emph{structurally inert}: the
profile inputs were populated only during training rollouts and were
zero in every evaluation, confirmation, and deployment path. The
network had been fitting a signal it never sees at test time.

\subsection{Search in training does not pay}

Each lineage's final feature configuration was trained a further $120$k
games with IS-MCTS in the rollout loop (the learner's own seats
consult the tree below 26 cards; opponents play raw policies), then
gated against its own predecessor under the paired protocol. All six
paired seed-level differences were negative:

\begin{center}
\begin{tabular}{lccc c}
\toprule
& seed 11 & seed 22 & seed 33 & mean \\
\midrule
A: search $-$ no-search & $-0.083$ & $-0.377$ & $-0.403$ & $-0.288$ \\
B: search $-$ no-search & $-0.029$ & $-0.192$ & $-0.239$ & $-0.154$ \\
\bottomrule
\end{tabular}
\end{center}

The cost is also large: rollout throughput fell from $\sim 140$
games/second (policy-only, $d{=}192$, 20 workers) to $12$--$19$
games/second. We offer two hypotheses. First, PPO's importance ratio is
anchored to the policy that generated the batch; when the executed move
comes from the tree, the behaviour and target distributions separate,
and the clipped objective is being asked to learn from off-policy
actions it was not designed for. Second, and consistent with
\S\ref{sec:census}, the tree changes so few decisions that the signal
it injects is small relative to the variance it adds.

\subsection{Where search actually acts}
\label{sec:census}

We instrumented 300 games per configuration, recording every decision's
source. Table~\ref{tab:census} gives the breakdown by cards remaining.

\begin{table}[h]
\centering
\begin{tabular}{lrrrrr}
\toprule
Cards left & forced & policy & search override & solver & $n$ \\
\midrule
52--40 & 194 & 578 & --- & --- & 772 \\
39--27 & 379 & 674 & --- & --- & 1053 \\
26--23 & 167 & 183 & 5 & 0 & 355 \\
22--19 & 201 & 152 & 1 & 2 & 356 \\
18--15 & 207 & 145 & 2 & 4 & 358 \\
14--1 & 345 & 26 & 4 & 155 & 530 \\
\midrule
total & 43.6\% & 51.4\% & \textbf{0.4\%} & 4.7\% & 3424 \\
\bottomrule
\end{tabular}
\caption{Decision sources by cards remaining (search enabled below 26).
Rows 1--2 are outside the search band.}
\label{tab:census}
\end{table}

The override rate is remarkably stable: $0.4\%$ at 32 simulations with
a $0.35$ margin, $0.4\%$ at 128 simulations with a $0.10$ margin,
$0.4\%$ at 256 simulations, and $0.4\%$ at depth 16 with the
trick-boundary cutoff disabled --- a $32\times$ compute range and a
$3.5\times$ margin range with no measurable effect. Against a weak,
out-of-distribution room (scripted heuristics, a linear champion, and a
random player) the rate doubles to $0.9\%$, and the character of the
overrides changes: substitutions of one single for another give way to
\emph{single $\rightarrow$ pass} corrections.

The exact solver, by contrast, is decisive: $155$ of the $167$ active
(non-forced, non-policy) interventions occur at $\leq 14$ cards, where
it makes $81\%$ of all non-forced decisions, and roughly half of its
moves are \emph{boss} plays --- provably unbeatable singles, pairs, or
five-card hands. We read this as: \emph{the value of "search" in this
domain is overwhelmingly the value of exact endgame solving}, and the
sampled tree above it functions mainly as a verifier of an already
well-trained policy.

\subsection{Non-transitivity}

We ran all $\binom{4+3}{4} = 35$ multisets of four agents over the four
seats, $1000$ games each with seats reshuffled per game, plus reference
matches against triples of the two shipped predecessors. The agents are
A (full stack), B (lean/exploiter lineage), C (a humanlike-anchored
agent trained with an imitation loss toward a clone of strong human
players), and D (trained explicitly against A, B, and C at one table).

The pairwise structure is cyclic: B outscores A ($+0.40$ vs $-0.13$
with three A seats and one B), A dominates C ($+0.95$ vs $-0.32$),
and C punishes B ($+0.68$ vs $-0.15$). Composition matters
independently of pairwise strength: B, the strongest agent
head-to-head against both shipped predecessors ($+0.94$ and $+0.59$),
is the \emph{worst} performer at the all-different table
($-0.51$). Agent D, trained specifically to beat the other three,
achieved only small edges ($+0.17$ to $+0.27$) --- confirming the
project's prior that a well-trained field is hard to exploit --- but
paid for its specialization with $-0.66$ against a mirror room of the
older champion.

\subsection{Objective blunders}

We measure three solver-verified error classes that require no
judgment: playing something other than an available guaranteed win;
gifting a beatable move to a player one card from out while holding an
unanswerable alternative; and spending a beatable move while holding a
boss unit with an opponent two cards from out. Over $\sim 1650$
decisions per agent, all three shipped agents record \textbf{0.00
missed guaranteed wins} per 100 decisions and $0.00$--$0.36$ gifts, but
$1.5$--$2.2$ wasted-boss errors. The endgame is solved; control timing
is not.

\subsection{Human study}

Against 570 recorded games from human players, we replayed every
decision where a human diverged from the agent, played both choices out
over belief-sampled deals, and kept those where the human's move
measurably won. Humans diverged at 378 decisions and were measurably
better at 163 ($43\%$). The two dominant corrections were \emph{pass
instead of playing a single} (58 cases) and \emph{play a lower single}
(47 cases; the humans' card averaged $0.66$ ranks lower). By stage:
76 early ($>36$ cards), 64 middle, and only 23 late ($\leq 20$). The
residual human edge is early-hand patience and rank conservation ---
precisely the region where the exact solver cannot reach and where, per
\S\ref{sec:census}, the tree almost never overrides.

\section{Discussion}

The picture that emerges is that in a four-player climbing game with a
strong learned prior, the components divide cleanly by horizon.
\emph{Exact solving} owns the endgame and is worth every cycle:
zero missed guaranteed wins, half of its moves provably unbeatable.
\emph{The learned policy} owns the opening and midgame, and once its
training diet is diverse enough, a sampled tree at 250\,ms
per move can barely improve on it in-distribution. \emph{The tree} is
best understood not as a strength multiplier but as
\emph{insurance}: its override rate doubles and its corrections become
qualitatively different (patience rather than substitution) exactly
when the opponents leave the training distribution --- which is the
situation a deployed agent faces against humans.

This suggests a deployment policy that differs from the training
policy, and our final configuration reflects it: policy-only training
(fast, stable), then search at every decision at play time, where the
opponents are unknown and a second opinion is cheap.

Two methodological lessons generalize beyond Big~2. First, in
four-player games the evaluation seed determines table composition, and
composition effects are large enough to reverse conclusions; paired
multi-seed evaluation is not optional. Second, engineered features must
be validated in the pipeline that will actually run them --- our
opponent-profile block trained on inputs that were identically zero at
evaluation time, an error invisible to training curves.

\section{Future work: bridge and mahjong}

Two natural next domains stress different parts of this design.

\paragraph{Bridge} adds a partnership and an explicit communication
channel. The bidding phase is a cheap-talk game where the information
content of a bid is conventional rather than physical, so the
belief head would need to model \emph{partner's model of us}, not just
card distributions. Our finding that search functions as
out-of-distribution insurance is especially relevant: against
unfamiliar conventions, a determinized search whose world sampler
respects the auction is precisely the safety net that a policy trained
on a fixed convention set would lack. The exact-solving result should
transfer directly --- double-dummy solvers are mature, and the
analogous claim would be that a policy plus double-dummy endgame beats
a policy alone by more than a sampled tree adds in the middle game.

\paragraph{Mahjong} preserves the four-player non-partnership structure
that motivated this study, while replacing the climbing dynamic with
set collection under a wall of hidden tiles and interrupt-driven turn
order (pon/chi/kan). Suphx \citep{li2020} established that deep RL with
oracle distillation is competitive; what our results suggest is worth
testing there is (i) whether the same policy--search convergence
appears once the opponent diet is sufficiently diverse, (ii) whether
mahjong's much larger hidden-state space pushes the crossover point
where sampled search stops agreeing with the policy, and (iii) whether
an exactly solvable endgame region exists at all --- if it does not,
our conclusion that ``search value is mostly solver value'' would
predict a substantially different component balance. Mahjong's
per-seat scoring cliffs (yakuman, riichi sticks) also make the tiered
payout structure we studied here a natural precursor.

\section{Limitations}

Our evaluation pool is largely self-referential: agents are scored
against other agents from the same project, and absolute scores are not
comparable to published Big~2 systems. The human study covers 570 games
from a small number of players. All measurements come from a single
hardware configuration, and search-cost conclusions are therefore
compute-relative. Finally, the training runs reported here are short by
the standards of the field ($10^5$ games, not $10^7$); the negative
result on search-in-training may not hold at scale, where the policy
has more capacity to absorb the tree's corrections.

\paragraph{Reproducibility.} All code, trained checkpoints, evaluation
harnesses, and the tournament and census scripts are available at
\url{https://github.com/bluo42/big2-ai}.

\begin{thebibliography}{9}
\bibitem[Frank \& Basin(1998)]{frank1998} I.~Frank and D.~Basin.
  Search in games with incomplete information. \emph{Artificial
  Intelligence}, 1998.
\bibitem[Long et al.(2010)]{long2010} J.~Long, N.~Sturtevant,
  M.~Buro, and T.~Furtak. Understanding the success of perfect
  information Monte Carlo sampling. \emph{AAAI}, 2010.
\bibitem[Zha et al.(2021)]{zha2021} D.~Zha et al. DouZero: Mastering
  DouDizhu with self-play deep reinforcement learning. \emph{ICML},
  2021.
\bibitem[Yang et al.(2022)]{yang2022} G.~Yang et al. PerfectDou:
  Dominating DouDizhu with perfect information distillation.
  \emph{NeurIPS}, 2022.
\bibitem[Li et al.(2020)]{li2020} J.~Li et al. Suphx: Mastering
  Mahjong with deep reinforcement learning. arXiv:2003.13590, 2020.
\bibitem[Brown \& Sandholm(2019)]{brown2019} N.~Brown and
  T.~Sandholm. Superhuman AI for multiplayer poker. \emph{Science},
  2019.
\bibitem[Schmid et al.(2021)]{schmid2021} M.~Schmid et al. Player of
  Games. arXiv:2112.03178, 2021.
\bibitem[Charlesworth(2018)]{charlesworth2018} H.~Charlesworth.
  Application of self-play reinforcement learning to a four-player
  game of imperfect information. arXiv:1808.10442, 2018.
\end{thebibliography}

\end{document}
